\section{Diseño e Implementación del Sistema}
\label{sec:implementacion}

\subsection{Arquitectura Modular del Sistema (4 Capas)}
El sistema computacional ha sido concebido bajo una arquitectura desacoplada en cuatro capas con dependencias estrictamente unidireccionales y descendentes (Figura~\ref{fig:arquitectura}). El núcleo de cálculo de caminos mínimos es completamente independiente de la interfaz gráfica web, de los mecanismos de almacenamiento en disco y de la estructura concreta de cola de prioridad empleada. La inyección abstracta de colas permite evaluar y sustituir componentes sin modificar el algoritmo de Dijkstra ni reescribir los módulos de ruteo.

\begin{figure}[H]
    \centering
    \begin{tikzpicture}[
        layer/.style={rectangle, very thick, minimum width=15.4cm, minimum height=1.85cm, rounded corners=6pt},
        box/.style={rectangle, draw=black!60, fill=white, thick, minimum height=0.85cm, text width=4.35cm, align=center, rounded corners=4pt, inner ysep=4pt, inner xsep=4pt},
        arrow/.style={-{Latex[length=3.2mm, width=2.0mm]}, line width=1.2pt, draw=black!75}
    ]
    \node[layer, fill=purple!5, draw=purple!70] (L4) at (0, 8.20) {};
    \node[anchor=north, yshift=-4pt, text=purple!85!black, font=\scriptsize\bfseries] at (L4.north) {4. VALIDACIÓN, EXPERIMENTACIÓN Y APLICACIÓN};
    \node[box] at (-5.0, 8.02) {{\fontsize{8.3}{10}\selectfont\textbf{\texttt{tests/}}}\\[1pt]{\fontsize{7.3}{9}\selectfont 278 pruebas automatizadas}};
    \node[box] at (0, 8.02) {{\fontsize{8.3}{10}\selectfont\textbf{\texttt{scalability\_benchmark.py}}}\\[1pt]{\fontsize{7.3}{9}\selectfont Benchmark Canónico (1350 corridas)}};
    \node[box] at (5.0, 8.02) {{\fontsize{8.3}{10}\selectfont\textbf{\texttt{app.py}}}\\[1pt]{\fontsize{7.3}{9}\selectfont Streamlit Web + Plotly}};

    \node[layer, fill=green!5, draw=green!60!black] (L3) at (0, 5.55) {};
    \node[anchor=north, yshift=-4pt, text=green!50!black, font=\scriptsize\bfseries] at (L3.north) {3. ALGORITMOS Y VERIFICACIÓN FORMAL};
    \node[box] at (-5.0, 5.37) {{\fontsize{8.3}{10}\selectfont\textbf{\texttt{dijkstra.py}}}\\[1pt]{\fontsize{7.3}{9}\selectfont SSSP Parametrizado}};
    \node[box] at (0, 5.37) {{\fontsize{8.3}{10}\selectfont\textbf{\texttt{a\_star.py}}}\\[1pt]{\fontsize{7.3}{9}\selectfont $A^*$ con Heurística Consistente}};
    \node[box] at (5.0, 5.37) {{\fontsize{8.3}{10}\selectfont\textbf{\texttt{certificate.py}}}\\[1pt]{\fontsize{7.3}{9}\selectfont Certificado Bellman $O(|V|+|E|)$}};

    \node[layer, fill=orange!6, draw=orange!75!black] (L2) at (0, 2.90) {};
    \node[anchor=north, yshift=-4pt, text=orange!85!black, font=\scriptsize\bfseries] at (L2.north) {2. ESTRUCTURAS DE DATOS Y COLAS DE PRIORIDAD};
    \node[box] at (-5.0, 2.72) {{\fontsize{8.3}{10}\selectfont\textbf{\texttt{IndexedMinHeap}}}\\[1pt]{\fontsize{7.3}{9}\selectfont Min-Heap Binario Indexado}};
    \node[box] at (0, 2.72) {{\fontsize{8.3}{10}\selectfont\textbf{\texttt{DAryHeap} ($d=4$)}}\\[1pt]{\fontsize{7.3}{9}\selectfont Montículo 4-ario en Caché}};
    \node[box] at (5.0, 2.72) {{\fontsize{8.3}{10}\selectfont\textbf{\texttt{FibonacciHeap}}}\\[1pt]{\fontsize{7.3}{9}\selectfont Cortes en Cascada $O(1)^*$}};

    \node[layer, fill=blue!5, draw=blue!70!black] (L1) at (0, 0.25) {};
    \node[anchor=north, yshift=-4pt, text=blue!85!black, font=\scriptsize\bfseries] at (L1.north) {1. DATOS TOPOGRÁFICOS, COSTOS Y SERVICIOS};
    \node[box] at (-5.0, 0.07) {{\fontsize{8.3}{10}\selectfont\textbf{\texttt{GraphInstance} / \texttt{cost\_model}}}\\[1pt]{\fontsize{7.3}{9}\selectfont Malla Vial e Impedancia}};
    \node[box] at (0, 0.07) {{\fontsize{8.3}{10}\selectfont\textbf{\texttt{osm\_loader.py}}}\\[1pt]{\fontsize{7.3}{9}\selectfont OSMnx + Serialización JSON}};
    \node[box] at (5.0, 0.07) {{\fontsize{8.3}{10}\selectfont\textbf{\texttt{weather.py}}}\\[1pt]{\fontsize{7.3}{9}\selectfont Doble Proveedor + Offline Snapshot}};

    \draw[arrow] (L4.south) -- (L3.north) node[midway, right=8pt, font=\scriptsize\itshape] {ejecuta e instrumenta};
    \draw[arrow] (L3.south) -- (L2.north) node[midway, right=8pt, font=\scriptsize\itshape] {invoca operaciones de cola};
    \draw[arrow] (L2.south) -- (L1.north) node[midway, right=8pt, font=\scriptsize\itshape] {consume pesos y topología};
    \end{tikzpicture}
    \caption{Arquitectura modular desacoplada del sistema en cuatro capas operativas}
    \label{fig:arquitectura}
\end{figure}

\subsection{Descripción Detallada de Módulos y Componentes}
El software se estructura modularmente dentro del directorio \texttt{python/}:
\begin{enumerate}
    \item \texttt{src/algorithms/dijkstra.py}: Motor algorítmico principal que implementa el algoritmo de Dijkstra con inyección de colas de prioridad y telemetría de operaciones elementales.
    \item \texttt{src/algorithms/a\_star.py}: Búsqueda heurística $A^*$ con precomputación del factor admisible $\lambda$, construcción de heurística consistente euclidiana escalada y métricas de poda.
    \item \texttt{src/algorithms/reference.py}: Oráculos de referencia independientes (Bellman-Ford y Dijkstra ingenuo sobre lista no ordenada).
    \item \texttt{src/datastructures/}: Implementación pura de \texttt{IndexedMinHeap}, \texttt{DAryHeap} ($d=4$) y \texttt{FibonacciHeap} bajo contrato polimórfico común.
    \item \texttt{src/io/graph\_io.py}: Ingestión, parsing y validación de grafos en JSON, listas de adyacencia y formato de red vial andina.
    \item \texttt{src/io/osm\_loader.py}: Ingestión, reproyección y serialización de redes OpenStreetMap mediante OSMnx, con soporte nativo de \textit{snapshots} desconectados.
    \item \texttt{src/services/weather.py}: Servicio meteorológico desacoplado con arquitectura de doble proveedor (Open-Meteo REST API y OpenWeatherMap API v2.5), persistencia en \textit{snapshots} y respaldo ante fallos de conexión.
    \item \texttt{src/services/traffic.py}: Modelo de congestión vehicular horaria y coeficientes por clase de vía.
    \item \texttt{src/verification/certificate.py}: Validador del certificado formal de optimalidad en $O(|V|+|E|)$ y clasificador de equivalencia de soluciones.
    \item \texttt{src/benchmarks/scalability\_benchmark.py}: Banco de experimentación sistemática para medir tiempos y primitivas en perfiles sintéticos dispersos y densos.
    \item \texttt{app.py}: Aplicación web interactiva desarrollada en Streamlit con visualizaciones geográficas Plotly y cinco módulos operativos.
    \item \texttt{main.py} y \texttt{src/cli.py}: Interfaz de línea de comandos para automatización y demostración en vivo.
\end{enumerate}

\subsection{Formatos y Estructura de Datos (Serialización JSON)}
Todos los artefactos de datos se serializan en formato JSON estructurado y versionado, garantizando que las pruebas no dependan de APIs externas en vivo (Tabla~\ref{tab:formatos}).

\begin{table}[H]
    \centering
    \fontsize{9}{11}\selectfont
    \caption{Archivos de datos principales y su rol funcional en el sistema}
    \label{tab:formatos}
    \begin{tabularx}{\textwidth}{@{}>{\ttfamily\raggedright\arraybackslash}p{5.1cm} >{\raggedright\arraybackslash}X@{}}
        \toprule
        \textbf{Archivo en Repositorio} & \textbf{Contenido y Propósito} \\
        \midrule
        \url{data/instances/cusco_centro.json} & Instancia nuclear didáctica (6 nodos, 9 aristas, longitudes, pendientes y tiempos). \\
        \url{data/processed/cusco_osm_drive.json} & Snapshot oficial de Cusco OSM (891 nodos con lat/lon WGS84, 2,139 aristas dirigidas). \\
        \url{data/processed/benchmark_results.json} & Protocolo experimental, ambiente, muestras cronometradas, medianas e IQR de las 10 escalas. \\
        \url{data/processed/applied_results.json} & Mediciones de la consulta OSM, comparación de colas, poda de $A^*$ y sensibilidad climática. \\
        \url{data/processed/weather_cusco.json} & Snapshot meteorológico local actualizado de contingencia. \\
        \bottomrule
    \end{tabularx}
\end{table}

\begin{lstlisting}[language=Python, caption={Extracto del esquema JSON serializado para un segmento vial cusqueño}]
{
  "origin": "A",
  "destination": "B",
  "length_m": 180.0,
  "street_type": "calle_centro",
  "gradient_pct": 1.5,
  "street_name": "Calle Espaderos / Portal Carnicerias",
  "travel_time_sec": 44.82
}
\end{lstlisting}

\subsection{Contratos de Interfaz e Inyección de Colas de Prioridad}
Para posibilitar la sustitución dinámica y transparente de colas en el solver de Dijkstra sin alterar el flujo algorítmico, cada una de las tres estructuras implementa una interfaz uniforme:
\begin{itemize}
    \item \texttt{insert(item: str, key: float) -> None}: Inserta un identificador de vértice con su distancia tentativa asociada.
    \item \texttt{extract\_min() -> Tuple[float, str]}: Extrae y retorna el par $(\text{clave}, \text{elemento})$ de menor valor acumulado, reestructurando el montículo según sus invariantes.
    \item \texttt{decrease\_key(item: str, new\_key: float) -> None}: Disminuye la clave del vértice a su nuevo valor relajado, flotando el nodo o ejecutando cortes en cascada.
    \item \texttt{is\_empty() -> bool}: Determina en $O(1)$ si existen vértices pendientes de procesamiento.
    \item \texttt{op\_stats() -> Dict[str, int]}: Retorna la telemetría acumulada de inserciones, extracciones, decrementos de clave y comparaciones.
\end{itemize}

\subsection{Manejo Defensivo de Entradas y Validación}
La validación centralizada en \texttt{src/io/graph\_io.py} protege al sistema contra datos corruptos, rechazando de forma determinista: JSON mal formado, campos requeridos faltantes, identificadores duplicados, pesos no numéricos, NaN o infinitos, costos negativos (en Dijkstra) y nodos de origen/destino inexistentes. La CLI retorna códigos de salida estandarizados: 0 (éxito), 1 (fallo en la verificación del certificado de optimalidad) y 2 (entrada inválida o error de sintaxis).

\subsection{Servicios Complementarios: Meteorología y Escenarios de Tráfico}
\label{sec:servicio_meteorologico}
El módulo \texttt{src/services/weather.py} implementa un desacoplamiento mediante una arquitectura de doble proveedor:
\begin{enumerate}
    \item \textbf{Open-Meteo REST API \parencite{openmeteo2026}:} Fuente pública libre sin autenticación basada en datos de modelos meteorológicos de predicción numérica.
    \item \textbf{OpenWeatherMap API (v2.5 Current Weather) \parencite{openweathermap2026}:} Proveedor autenticado mediante clave de API suministrada exclusivamente por variables de entorno locales, con mapeo hacia \texttt{WeatherObservation}.
\end{enumerate}
Ambos clientes alimentan la función multiplicadora $\Omega_{\text{clima}} \ge 1.0$:
\begin{equation}
    \Omega_{\text{clima}} = 1.0 + c_{\text{lluvia}} \min\left(\frac{p}{10}, 1.0\right) + c_{\text{vis}} \max\left(0, 1 - \frac{v}{10\,000}\right) + c_{\text{ráfaga}} \min\left(\frac{r}{80}, 1.0\right) + c_{\text{hum}} f(\text{HR})
\end{equation}
Si la conectividad a Internet falla o la clave presenta latencia de activación, el sistema opera sin interrupciones consumiendo los datos locales del \textit{snapshot}.

\subsection{Decisiones de Implementación Relevantes y Justificación Técnica}
\begin{itemize}
    \item \textbf{Programación pura sin \texttt{heapq}:} Las tres colas se programaron desde cero en Python puro para garantizar que las mediciones reflejen las propiedades estructurales de los algoritmos y no optimizaciones escritas en lenguaje C.
    \item \textbf{Degeneración admisible de $A^*$:} Si un grafo no posee coordenadas geográficas, $A^*$ adopta $h(v)=0$, degenerando de manera admisible en el algoritmo canónico de Dijkstra sin romper la ejecución.
    \item \textbf{Reproducibilidad sin dependencias en vivo:} La red vial urbana y los datos climáticos se congelan en archivos serializados en \texttt{data/processed/}, garantizando ejecuciones deterministas en cualquier computadora sin conexión a red.
    \item \textbf{Separación estricta entre núcleo y visualización:} La capa matemática y algorítmica es 100\% agnóstica de Streamlit, asegurando que la suite de pruebas unitarias y los benchmarks se ejecuten sin sobrecargas de renderizado web.
\end{itemize}
